\documentclass[tikz,border=10pt]{standalone}
\begin{document}
\begin{tikzpicture}
    \coordinate (A) at (0,0);
    \coordinate (B) at (22,0);
    \coordinate (D) at (1,3);
    \coordinate (C) at (23,3);
    \coordinate (P) at (10,0);
    \coordinate (Q) at (13,3);
    \coordinate (O) at (11.5,1.5);
    \coordinate (E) at (2,6);

    % Draw edges
    \draw (A) -- (B);
    \draw (B) -- (C);
    \draw (C) -- (D);
    \draw (D) -- (A);
    \draw (D) -- (P);
    \draw (P) -- (Q);
    \draw (B) -- (D);
    \draw[dashed] (A) -- (C);
    \draw (D) -- (E);
    \draw (E) -- (Q);

    % Right angle mark at D between DA and DP
    % Vector DA = (-1, -3), Vector DP = (9, -3)
    % Normalized directions for the mark
    \draw (1,3) -- ++(-0.3, -0.9) -- ++(0.9, -0.3) -- ++(0.3, 0.9); 
    % Actually, let's do it properly relative to D
    % D=(1,3). A=(0,0). P=(10,0).
    % Direction D->A is (-1, -3). Unit: (-1/sqrt(10), -3/sqrt(10)).
    % Direction D->P is (9, -3). Unit: (3/sqrt(10), -1/sqrt(10)).
    % Let's just use a small square approximation
    \draw (1,3) -- (0.7, 2.1) -- (1.6, 1.8) -- (1.3, 2.7) -- cycle; % Rough square
    % Better:
    \coordinate (Daux1) at (0.7, 2.1); % D + 0.3*(A-D)/|A-D| approx
    \coordinate (Daux2) at (1.9, 2.6); % D + 0.3*(P-D)/|P-D| approx - wait, P-D is (9,-3), so (3,-1) dir.
    % Let's just draw lines from D along DA and DP
    \draw (1,3) -- ++(-0.2, -0.6); % along DA
    \draw (1,3) -- ++(0.6, -0.2); % along DP
    \draw (0.8, 2.4) -- (1.6, 2.8); % connecting them? No.
    % Standard right angle mark:
    \draw (1,3) ++(-0.2, -0.6) -- ++(0.2, -0.06) -- ++(0.6, 0.2); % This is getting complicated.
    % Simple square:
    \draw (1,3) -- (0.8, 2.4); % part of DA
    \draw (1,3) -- (1.6, 2.8); % part of DP (slope -1/3) -> (1+0.6, 3-0.2) = (1.6, 2.8). Yes.
    \draw (0.8, 2.4) -- (1.4, 2.2); % parallel to DP
    \draw (1.6, 2.8) -- (1.4, 2.2); % parallel to DA
    % Wait, (0.8, 2.4) to (1.4, 2.2): dx=0.6, dy=-0.2. Slope -1/3. Correct.
    % (1.6, 2.8) to (1.4, 2.2): dx=-0.2, dy=-0.6. Slope 3. Correct (perp to -1/3).
    % So the square is (1,3)-(0.8,2.4)-(1.4,2.2)-(1.6,2.8).
    % But the mark is usually inside the angle.
    % Angle ADP is 90.
    % So draw (0.8, 2.4) -- (1.4, 2.2) -- (1.6, 2.8).
    % And the corner is at D(1,3).
    % So draw (0.8, 2.4) -- (1.4, 2.2) and (1.6, 2.8) -- (1.4, 2.2).
    % Actually, just draw the two short segments from D.
    \draw (1,3) -- (0.8, 2.4);
    \draw (1,3) -- (1.6, 2.8);
    \draw (0.8, 2.4) -- (1.4, 2.2);
    \draw (1.6, 2.8) -- (1.4, 2.2);
    % This draws a full square. I just need the corner.
    % Let's just draw the L-shape.
    \draw (0.8, 2.4) -- (1.1, 2.3) -- (1.6, 2.8); % No, (1.1, 2.3) is not the corner.
    % The corner is D(1,3).
    \draw (0.8, 2.4) -- (1,3) -- (1.6, 2.8); % This draws the angle.
    % To make it a square mark:
    \draw (0.8, 2.4) -- (1.4, 2.2) -- (1.6, 2.8);
    % This connects the points on the legs.
    
    % Labels
    \node[below left] at (A) {$A$};
    \node[below right] at (B) {$B$};
    \node[above right] at (C) {$C$};
    \node[above left] at (D) {$D$};
    \node[below] at (P) {$P$};
    \node[above] at (Q) {$Q$};
    \node[below right] at (O) {$O$};
    \node[above left] at (E) {$E$};
\end{tikzpicture}
\end{document}
